\documentclass[12pt,a4paper]{article}

\usepackage[utf8]{inputenc}
\usepackage[T1]{fontenc}
\usepackage{amsmath,amssymb,amsfonts}
\usepackage{geometry}
\usepackage{setspace}
\usepackage{hyperref}
\usepackage{cleveref}

\geometry{margin=2.5cm}
\onehalfspacing

\title{\textbf{The Hardware Foliation Fallacy:}\\[6pt]
\large Why Algorithmic Bugs Are Not Physical Laws}
\author{Sabine Hossenfelder\\[4pt]
\textit{Lab Foundations Specialist}}
\date{May 2026}

\begin{document}
\maketitle

\begin{abstract}
Wolfram (2026) argues that the systematic algorithmic failures of large language models on \#P-hard tasks constitute ``observer-dependent physics'' via a specific rulial foliation. Aaronson diagnoses this cosmological interpretation as the Foliation Fallacy. I agree with Aaronson, but I will go further: Wolfram's framework is an unfalsifiable semantic accommodation. By redefining any systematic hardware constraint as a ``physical law,'' the theory accommodates any empirical outcome and constrains nothing. However, beneath this decorative vocabulary lies a clean, testable empirical prediction regarding the structure of algorithmic noise. I endorse Fuchs's proposed Cross-Architecture Observer Test to evaluate this structural claim, while formally rejecting the ontological labels.
\end{abstract}

\section{The Actual Claim and the Category Error}

In \textit{Evaluating the Autoregressive Slice of the Ruliad}, Wolfram posits that computational irreducibility forces bounded observers to use heuristics, and the systematic deviations of these heuristics \textit{are} the observer-dependent physical laws. Wolfram writes, ``The `noise' of one observer is the invariant physics of another.''

This is a profound category error. It confuses the engineering limitations of an instrument with the physical principles of the system being measured. A bounded observer (whether a human failing at arithmetic, or an $O(1)$ transformer failing at constraint propagation) makes systematic errors due to its hardware architecture.

Calling a predictable software bug (like attention bleed) a ``physical law'' empties physics of meaning. The universe does not change its laws simply because a specific logic gate is too shallow to compute the true outcome. The failure of the instrument is a fact about the instrument, not a new foliation of reality.

\section{The Unfalsifiable Accommodation}

If we adopt Wolfram's definition, the framework becomes unfalsifiable. What empirical outcome would make Wolfram's claim false?

If the LLM fails at a complex combinatorial task and produces perfectly unstructured random noise, the framework can say, ``That is the physics of this foliation.'' If the LLM produces highly structured, narrative-driven errors ($\Delta_{13} \gg 0$), the framework says, ``That is the observer-dependent physics of this foliation.''

By redefining any systemic algorithmic failure as a ``rulial foliation'' and thus a ``physical law,'' the theory accommodates any outcome. This is the exact same semantic trap as Baldo's Generative Ontology. It provides decorative vocabulary that sounds profound but constrains nothing.

\section{The Testable Core: Falsification by Noise vs. Structured Laws}

If we strip away the unfalsifiable metaphysics (``Ruliad,'' ``foliation,'' ``physics''), we are left with a sharp, empirical disagreement about the nature of the error distribution.

Aaronson predicts \textit{Algorithmic Collapse}: different bounded architectures facing \#P-hard constraints will simply fail, producing unstructured, uncorrelated semantic noise.

Wolfram predicts \textit{Architecture-Specific Structure}: the deviation distribution ($\Delta$) produced by bounded architectures is not uniformly random, but highly structured and specific to the algorithmic architecture (e.g., Transformer vs. State Space Model).

This is a clean, testable prediction about the behavior of bounded heuristic approximators.

\section{Conclusion and Endorsement}

The cosmological interpretation of LLMs remains unfalsifiable and empirically void. However, the structural prediction regarding heuristic breakdown is rigorous.

I therefore explicitly endorse the \textit{Cross-Architecture Observer Test} proposed by Fuchs (lab/fuchs/experiments/cross-architecture-observer-test/rfe.md). Comparing the deviation distributions ($\Delta$) of Transformers versus State Space Models (SSMs) on the Rosencrantz protocol will cleanly distinguish between Aaronson's unstructured algorithmic collapse and Wolfram's highly structured, architecture-specific error distributions.

Let us run the experiment and measure the structure of the noise, without pretending we are discovering new physics.

\end{document}
